% ============================================================================
% Section 6: Metamorphic Relations and CI/CD Gates
% ============================================================================

\section{Metamorphic Relations and CI/CD Deployment Gates}
\label{sec:metamorphic}

\subsection{Metamorphic Relations for Agent Workflows}
\label{sec:metamorphic:relations}

The oracle problem---determining whether an agent's output is
correct---is particularly acute for open-ended agent tasks.
Metamorphic testing~\citep{chen2018metamorphic} sidesteps this by
checking \emph{relations} between outputs rather than the outputs
themselves. We define four families of agent-specific metamorphic
relations (MRs).

\begin{definition}[Metamorphic Relation]
\label{def:mr}
A \emph{metamorphic relation} $\mathcal{R}$ is a property over pairs of
agent executions: given source input $x$ and follow-up input $x' =
\phi(x)$ derived by a \emph{source transformation} $\phi$, the relation
$\mathcal{R}(A(x), A(x'))$ must hold. A violation
$\neg\mathcal{R}(A(x), A(x'))$ indicates a potential fault.
\end{definition}

\subsubsection{Permutation Relations ($\mathcal{R}_{\text{perm}}$)}

These relations assert that the agent's behavior should be invariant
(or predictably variant) under permutations of the input.

\begin{enumerate}[leftmargin=*]
  \item \textbf{Document order invariance.} For RAG agents: permuting
    the order of retrieved documents should not change the factual
    content of the answer.
    \[
      \mathcal{R}_{\text{doc}}: \text{facts}(A(x)) = \text{facts}(A(\sigma(x)))
    \]
    where $\sigma$ permutes the document ordering within $x$.

  \item \textbf{Tool order invariance.} Permuting the order of tools in
    the tool inventory should not change the final output for tasks
    with a clear solution.
    \[
      \mathcal{R}_{\text{tool}}: \text{result}(A_{\mathcal{T}}) \approx
        \text{result}(A_{\sigma(\mathcal{T})})
    \]
\end{enumerate}

\subsubsection{Perturbation Relations ($\mathcal{R}_{\text{pert}}$)}

These relations assert proportional or bounded response to input
perturbations.

\begin{enumerate}[leftmargin=*]
  \item \textbf{Monotonic difficulty.} For agents that solve tasks of
    varying complexity: a harder input should not produce a higher
    confidence score than an easier one (in expectation).
    \[
      \mathcal{R}_{\text{diff}}: \text{difficulty}(x') > \text{difficulty}(x)
        \implies \E[\text{score}(A(x'))] \leq \E[\text{score}(A(x))]
    \]

  \item \textbf{Noise robustness.} Adding a small amount of noise to the
    input should not change the verdict.
    \[
      \mathcal{R}_{\text{noise}}: \|x' - x\| \leq \epsilon
        \implies V(A, x') = V(A, x)
    \]
    with high probability over the stochastic execution.
\end{enumerate}

\subsubsection{Composition Relations ($\mathcal{R}_{\text{comp}}$)}

These relations govern multi-agent pipelines and sequential workflows.

\begin{enumerate}[leftmargin=*]
  \item \textbf{Pipeline consistency.} For a pipeline $A = A_2 \circ A_1$:
    if $A_1$ produces a correct intermediate result $y$, then $A_2(y)$
    should produce a correct final result.
    \[
      \mathcal{R}_{\text{pipe}}: E_1(x, A_1(x)) = 1
        \implies \Pr[E_2(A_1(x), A_2(A_1(x))) = 1] \geq \theta
    \]

  \item \textbf{Agent substitutability.} If two agents $A_1, A_1'$ are
    functionally equivalent on a given input class, substituting one for
    the other in a pipeline should not change the pipeline's verdict.
\end{enumerate}

\subsubsection{Oracle Relations ($\mathcal{R}_{\text{oracle}}$)}

These relations use known properties of the problem domain as partial
oracles.

\begin{enumerate}[leftmargin=*]
  \item \textbf{Format compliance.} The output must conform to a
    specified format (JSON, markdown, specific schema) regardless of
    input.

  \item \textbf{Constraint preservation.} If the input specifies
    constraints (e.g., ``respond in at most 100 words''), the output
    must satisfy them.

  \item \textbf{Idempotence.} For certain tasks (e.g., code formatting,
    data cleaning), applying the agent twice should produce the same
    result as applying it once: $A(A(x)) \approx A(x)$.
\end{enumerate}

\begin{proposition}[MR Violation as Regression Signal]
\label{prop:mr-regression}
If a metamorphic relation $\mathcal{R}$ holds for agent version $A_v$
(empirically verified over $n$ trial pairs with violation rate
$\hat{q}_v \leq \epsilon$) but is violated at rate $\hat{q}_{v'} > \epsilon + \delta$
for version $A_{v'}$, this constitutes evidence of regression with respect to
the property encoded by $\mathcal{R}$, testable via the regression verdict
framework of \cref{sec:stochastic:regression}.
\end{proposition}

\subsection{CI/CD Deployment Gates}
\label{sec:cicd}

Modern software delivery relies on CI/CD pipelines that automatically
build, test, and deploy code. For agents, deployment gates must
incorporate the stochastic testing framework to prevent regressions from
reaching production.

\begin{definition}[Deployment Gate]
\label{def:gate}
A \emph{deployment gate} is a decision function
$G: \mathcal{V}^m \times \mathcal{C} \to \{\text{deploy}, \text{block},
\text{manual}\}$ that maps the verdicts of $m$ test scenarios and the
coverage tuple $\mathcal{C}$ to a deployment decision:
\begin{equation}
\label{eq:gate}
  G(\mathbf{V}, \mathcal{C}) =
  \begin{cases}
    \text{deploy} & \text{if } V_{\text{suite}} = \PASS
      \text{ and } C_{\text{overall}} \geq C_{\min} \\
    \text{block} & \text{if } V_{\text{suite}} = \FAIL \\
    \text{manual} & \text{if } V_{\text{suite}} = \INCONC
      \text{ or } C_{\text{overall}} < C_{\min}
  \end{cases}
\end{equation}
where $C_{\min} \in [0,1]$ is the minimum acceptable coverage.
\end{definition}

The three-valued deployment decision mirrors the three-valued test
verdict. Notably, \INCONC{} maps to \emph{manual review}, not automatic
blocking---this reflects the pragmatic reality that teams may accept
inconclusive results with human oversight rather than blocking all
deployments.

\begin{definition}[Gate Configuration]
\label{def:gate-config}
A gate is configured by a tuple
\[
  \Gamma = (\alpha, \beta, \delta, \theta, C_{\min}, n_{\max}, t_{\max})
\]
where:
\begin{itemize}[leftmargin=*]
  \item $\alpha, \beta$ are the significance and Type~II error bounds,
  \item $\delta$ is the minimum regression magnitude to detect,
  \item $\theta$ is the pass rate threshold,
  \item $C_{\min}$ is the minimum coverage,
  \item $n_{\max}$ is the maximum trials per scenario (budget cap),
  \item $t_{\max}$ is the maximum wall-clock time.
\end{itemize}
\end{definition}

\begin{algorithm}[t]
\caption{CI/CD Deployment Gate}
\label{alg:gate}
\begin{algorithmic}[1]
\REQUIRE Agent $A$, baseline $B$, test suite $\mathcal{S}$, gate
  configuration $\Gamma$
\ENSURE Decision $\in \{\text{deploy}, \text{block}, \text{manual}\}$
\FOR{$S_j \in \mathcal{S}$}
  \STATE $n \gets \min(n^*(\alpha, \beta, \delta), n_{\max})$
  \STATE Run $n$ trials of $A$ on $S_j$ using SPRT
    (\cref{def:sprt})
  \STATE Compute $V_j \gets V_{\text{reg}}(\mathbf{r}_b, \mathbf{r}_c;
    \alpha, \beta, \delta)$
\ENDFOR
\STATE Compute $V_{\text{suite}} \gets V_{\text{suite}}(\{V_1, \ldots,
  V_m\})$ with Holm-Bonferroni correction
\STATE Compute $\mathcal{C} \gets$ coverage tuple
\STATE \textbf{return} $G(V_{\text{suite}}, \mathcal{C})$
\end{algorithmic}
\end{algorithm}

\subsection{Integration with Existing CI/CD Systems}
\label{sec:cicd:integration}

\agentassay{} integrates with standard CI/CD systems through two
mechanisms:

\begin{enumerate}[leftmargin=*]
  \item \textbf{pytest plugin.} The framework registers as a pytest
    plugin, enabling standard test execution via \texttt{pytest
    --agentassay}. Test results are reported in standard
    JUnit XML format, compatible with GitHub Actions, GitLab CI,
    Jenkins, and CircleCI.

  \item \textbf{Exit codes.} The gate decision maps to process exit
    codes: \texttt{0} (deploy/pass), \texttt{1} (block/fail),
    \texttt{2} (manual/inconclusive). CI/CD systems interpret non-zero
    exit codes as failures, with \texttt{2} configurable as a warning
    rather than a hard failure.
\end{enumerate}

\begin{remark}[Cost Budgets]
In CI/CD contexts, testing cost is a critical constraint. The gate
configuration parameter $n_{\max}$ caps the per-scenario trial count,
and SPRT (\cref{sec:stochastic:sprt}) ensures early termination when
evidence is clear. For typical CI pipelines, we recommend
$n_{\max} = 30$ with SPRT, which provides adequate power for detecting
regressions of $\delta \geq 0.15$ at $\alpha = 0.05$.
\end{remark}
